\chapter{Environmental and Resource Economics}\label{ch:b3:environmental-and-resource-economics}

The Year 2 volume set out the tools for an externality: a corrective tax, tradable
permits, bargaining when rights are clear, and the choice between prices and
quantities when costs are uncertain. This chapter takes them to the evidence. How
are damages valued? Did trading cut the cost of cleaning the air? Does pollution
fall as countries grow richer? What is a tonne of carbon dioxide worth, and why do
the estimates differ so much? And what does theory say about the price of a resource
that runs out --- and why does the record disagree? The physical science of the
climate is the Earth and space sciences books'; here it is taken as given.

\begin{openingchart}{Sulphur dioxide emitted per person against GDP per person (logarithmic
scale, 2011 dollars), six economies, every fifth year from 1850 (or the first year with
data) to 2019; an annual peak can fall between two points.}{Community Emissions Data System (CEDS); Maddison Project Database 2023}
\begin{tikzpicture}
\begin{axis}[omrecord, width=0.72\linewidth, height=0.46\linewidth, xmin=2.85,
  xmax=4.85, ymin=0, ymax=200, xtick={3,3.5,4,4.5}, xticklabels={1{,}000, 3{,}000, 10{,}000, 30{,}000},
  xlabel={GDP per person, dollars}, ylabel={SO$_2$, kg per person},
  legend pos=north west, legend style={font=\footnotesize}, legend columns=2]
  \addplot[black, thick, mark=*, mark size=0.8pt] table[x=lx, y=kg] {figdata/bachelor-3/35-environmental-and-resource-economics/out/so2_GBR.csv};
  \addplot[omDef, thick, mark=*, mark size=0.8pt] table[x=lx, y=kg] {figdata/bachelor-3/35-environmental-and-resource-economics/out/so2_USA.csv};
  \addplot[omDef, dashed, thick] table[x=lx, y=kg] {figdata/bachelor-3/35-environmental-and-resource-economics/out/so2_DEU.csv};
  \addplot[omThm, thick] table[x=lx, y=kg] {figdata/bachelor-3/35-environmental-and-resource-economics/out/so2_JPN.csv};
  \addplot[omThm, dashed, thick] table[x=lx, y=kg] {figdata/bachelor-3/35-environmental-and-resource-economics/out/so2_KOR.csv};
  \addplot[black!50, dotted, very thick] table[x=lx, y=kg] {figdata/bachelor-3/35-environmental-and-resource-economics/out/so2_CHN.csv};
  \legend{United Kingdom, United States, Germany, Japan, South Korea, China}
\end{axis}
\end{tikzpicture}
\end{openingchart}

% figdata: so2_stats
Each economy's emissions per person rose, peaked and fell. The United Kingdom peaked in 1907 at 134 kg per
person, with income at 7{,}855 dollars; the United States in 1918 at 191 kg; Germany in 1973 and Japan in 1969, at 107
and 51 kg; South Korea in 1979 and China in 2006, at incomes of 6{,}266 and 7{,}271 dollars. In 2019 all six were below 9 kg.
The peaks came at very different incomes.

\begin{recall}
Externalities are the grades 10--12 volume's. The corrective tax, Coase's theorem,
tradable permits and the comparison of prices and quantities are the Year 2
volume's; learning by doing is the Year 1 volume's. Discounting and the \omterm{def:b3:cost-benefit-analysis:vsl}{value of a
statistical life} are \cref{ch:b3:cost-benefit-analysis}'s.
\end{recall}

\begin{notation}[Symbols of this chapter]\label{not:b3:environmental-and-resource-economics:symbols}
$x_i$ is firm $i$'s abatement; $\mathrm{MAC}_i$ its \omterm{def:b3:environmental-and-resource-economics:abatement}{marginal abatement cost}; $\mathrm{MD}$ the marginal
damage; $\rho$ the rate of pure time preference; $\eta$ the elasticity of marginal utility;
$g$ the growth of consumption per person; $P_t$ a resource price; $c$ its extraction cost.
\end{notation}

\section{Valuing damages}\label{sec:b3:environmental-and-resource-economics:damages}

\begin{definition}[Damage function; marginal abatement cost]\label{def:b3:environmental-and-resource-economics:abatement}
\index{damage function}\index{marginal abatement cost}
A \emph{damage function} gives the money value of the harm done by a level of
emissions. A firm's \emph{marginal abatement cost} is the cost of cutting its
emissions by one more unit, given what it has already cut.
\end{definition}

\begin{measure}[How damages are valued]
Damages are built from three pieces: the change in exposure (a pollutant's
concentration where people live), a dose--response relation from epidemiology
(extra deaths or illnesses per unit of exposure), and a money value for each
outcome --- for deaths, the \omterm{def:b3:cost-benefit-analysis:vsl}{value of a statistical life} of
\cref{ch:b3:cost-benefit-analysis}. Hedonic house prices give a second route, for
damages that people can see and move away from. The first route is only as good as
its dose--response study; the second misses harms people do not know about.
The sulphur dioxide record below shows how much the first route can move an answer:
when the market was created, its benefits were expected to be about equal to its
costs; once epidemiology linked fine sulphate particles to deaths, the estimated
health benefits alone reached 50 to more than 100 billion dollars a year.
% ledger: b3er:ss.benefit
\end{measure}

\section{Taxes, permits and standards}\label{sec:b3:environmental-and-resource-economics:instruments}

\begin{definition}[Emission standard]\label{def:b3:environmental-and-resource-economics:standard}
\index{emission standard}
An \emph{emission standard} limits each source directly: a \emph{performance
standard} caps emissions per unit of output; a \emph{technology standard} requires a
given technique, such as a scrubber.
\end{definition}

\begin{proposition}[Equal marginal costs]\label{prop:b3:environmental-and-resource-economics:equal-mac}
A total cut $X$ shared among firms with increasing \omterm{def:b3:environmental-and-resource-economics:abatement}{marginal abatement costs} costs least when every
firm's \omterm{def:b3:environmental-and-resource-economics:abatement}{marginal abatement cost} is the same. A uniform tax $t$ achieves it, since each firm
abates until $\mathrm{MAC}_i = t$; so does a market in permits, at the permit price. A standard that requires
the same cut of every firm does not, unless their costs are the same.
\end{proposition}

\begin{proof}
Minimise $\sum_i C_i(x_i)$ subject to $\sum_i x_i = X$. The Lagrangian conditions are $C_i'(x_i) = \lambda$ for every $i$, and with convex $C_i$
they are sufficient. If $C_i'(x_i) < C_j'(x_j)$, moving a small amount of abatement from $j$ to $i$ keeps $X$ and lowers the
total cost. A firm facing a tax $t$ minimises $C_i(x_i) + t(e_i - x_i)$ and sets $C_i' = t$; a firm in a permit market sets $C_i'$ equal to
the permit price for the same reason. Equal cuts give $C_i'(X/n)$, which differ across $i$ unless the $C_i$ do.
\end{proof}

\begin{record}[The U.S. sulphur dioxide market]
The Clean Air Act Amendments of 1990 created a cap-and-trade market for sulphur dioxide from
power plants, to curb acid rain. It worked well, with surprises: its main benefits
were not the ones it was designed for; much of its low cost came from the earlier
deregulation of railways, which made low-sulphur coal cheaper to ship; and later
court decisions and rules led to the collapse of the market. % ledger: b3er:ss
From 1990 to 2004 power plants' emissions fell 36~\%, from 15.9 to 10.2 million tons,
while coal-fired generation rose 25~\%, and compliance was nearly complete. Against
rules that would have prescribed how each plant cut, the market's cost savings are
estimated at at least 15~\% and perhaps as much as 90~\%, depending on the rule it is
compared with. Annual benefits are put at 59 to 116 billion dollars against costs of
0.5 to 2 billion; more than 95~\% of the benefits are to health, from fewer fine
particles, not to the lakes and forests the law was written for.
% ledger: b3er:ss.cost, b3er:ss.benefit
\end{record}

\section{Pollution and income}\label{sec:b3:environmental-and-resource-economics:ekc}

\begin{definition}[Environmental Kuznets curve]\label{def:b3:environmental-and-resource-economics:ekc}
\index{environmental Kuznets curve}
An \emph{environmental Kuznets curve} is an inverted-U relation between a pollutant
and income per person: emissions first rise with income and then fall.
\end{definition}

\begin{model}{The environmental Kuznets curve}
\assumes As income grows, the economy first expands dirty industry, then demands
cleaner air, shifts to services and adopts cleaner techniques.
\predicts Local pollutants rise and then fall with income.
\tested Across urban air pollution and river quality, most indicators first worsen
and then improve with income, with turning points mostly below 8{,}000 dollars per person (Grossman and
Krueger 1994). % ledger: b3er:gk
Sulphur dioxide per person in the opening chart has the same shape.
\breaks The turning points differ widely: in the opening chart, emissions per person
peaked at about 6{,}300 dollars in Korea, 10{,}500 in the United States, 14{,}100 in Japan
and 19{,}100 in Germany; and the United Kingdom and China, which peaked at similar
incomes (7{,}900 and 7{,}300 dollars) a century apart, did so at levels nearly five times apart,
134 and 29 kilograms per person. Income alone sets neither the turning point nor the
height of the peak; laws and the move away from coal are the obvious candidates
(exercise 9). % figdata: so2_stats Carbon dioxide, whose damage falls on everyone and is
not felt locally, rises with income across countries (next chart).
\end{model}

\begin{omchart}{Carbon dioxide emitted per person against GDP per person, 2019, 152 countries
of over one million people (both logarithmic scales), with the least-squares line.}{Community
Emissions Data System (CEDS); Maddison Project Database 2023}
\begin{tikzpicture}
\begin{axis}[omscatter, width=0.6\linewidth, height=0.42\linewidth,
  xlabel={GDP per person, dollars}, ylabel={CO$_2$, tonnes per person},
  xtick={3,4,5}, xticklabels={1{,}000, 10{,}000, 100{,}000},
  ytick={-2,-1,0,1}, yticklabels={0.01, 0.1, 1, 10}]
  \addplot table[x=lx, y=ly] {figdata/bachelor-3/35-environmental-and-resource-economics/out/co2.csv};
  \addplot[black!60, thick, no marks, sharp plot] table[x=lx, y=ly] {figdata/bachelor-3/35-environmental-and-resource-economics/out/co2_fit.csv};
\end{axis}
\end{tikzpicture}
\end{omchart}
% figdata: co2_stats

Across 152 countries in 2019, emissions per person rise in proportion to income or a little faster:
the elasticity is 1.11. The median is 0.29 tonnes per person in the poorest fifth and 8.4 in the richest.
There is no hump.

\section{The social cost of carbon}\label{sec:b3:environmental-and-resource-economics:scc}

\begin{definition}[Social cost of carbon]\label{def:b3:environmental-and-resource-economics:scc}
\index{social cost of carbon}
The \emph{social cost of carbon} is the present value of the damages caused by one
more tonne of carbon dioxide emitted in a given year: $\mathrm{SCC}_0 = \sum_{t \ge 1} \mathrm{MD}_t/\prod_{s \le t}(1 + r_s)$, where $\mathrm{MD}_t$ is the extra
damage in year $t$ from the tonne and $r_s$ the discount rate.
\end{definition}

\begin{definition}[Ramsey formula (discounting)]\label{def:b3:environmental-and-resource-economics:ramsey-discount}
\index{Ramsey formula (discounting)}
The \emph{Ramsey formula} sets the consumption discount rate to $r = \rho + \eta g$, the rate of pure time
preference plus the elasticity of marginal utility times the growth of consumption
per person.
\end{definition}

\begin{proposition}[The discount rate from the Euler equation]\label{prop:b3:environmental-and-resource-economics:ramsey}
With utility $\sum_t e^{-\rho t} u(c_t)$ and $u(c) = c^{1 - \eta}/(1 - \eta)$, a dollar at date $t$ is worth $e^{-\rho t}(c_t/c_0)^{-\eta}$ dollars today; with consumption growing at $g$,
this is $e^{-(\rho + \eta g)t}$: the discount rate is $\rho + \eta g$.
\end{proposition}

\begin{proof}
A small transfer $\varepsilon$ from date 0 to date $t$ changes welfare by $\varepsilon[e^{-\rho t}u'(c_t) - u'(c_0)]$ per unit, so a unit at $t$ is worth $e^{-\rho t}u'(c_t)/u'(c_0)$ units at 0.
With $u'(c) = c^{-\eta}$ and $c_t = c_0e^{gt}$, $u'(c_t)/u'(c_0) = e^{-\eta gt}$, and the value is $e^{-(\rho + \eta g)t}$. This is the Euler equation of the Year 2
volume's Ramsey model read as a price.
\end{proof}

The two parameters do different work: $\rho$ says how much a later generation counts
because it is later, $\eta$ how much an extra dollar is worth to someone richer. With $\rho = 0.1~\%$, $\eta = 1$ and
$g = 1.3~\%$, $r = 1.4~\%$; with $\rho = 1.5~\%$ and $\eta = 2$, $r = 4.1~\%$. Because a tonne does harm for centuries, that difference multiplies
the \omterm{def:b3:environmental-and-resource-economics:scc}{social cost of carbon} several times over (exercise 5).

\begin{omchart}{The social cost of carbon dioxide by year of emission at three near-term
discount rates (2020 dollars per tonne; values rounded to two significant
figures).}{U.S. Environmental Protection Agency (2023), Report on the Social Cost of
Greenhouse Gases, Table ES.1}
\begin{tikzpicture}
\begin{axis}[omtime, width=0.62\linewidth, height=0.4\linewidth, xmin=2020,
  xmax=2080, ymin=0, ymax=650, ylabel={Dollars per tonne},
  legend pos=north west, legend style={font=\footnotesize}]
  \addplot[omThm, thick, mark=*, mark size=1.2pt] table[x=year, y=r15] {figdata/bachelor-3/35-environmental-and-resource-economics/out/scc.csv};
  \addplot[omDef, thick, mark=square*, mark size=1.2pt] table[x=year, y=r20] {figdata/bachelor-3/35-environmental-and-resource-economics/out/scc.csv};
  \addplot[black, thick, mark=triangle*, mark size=1.4pt] table[x=year, y=r25] {figdata/bachelor-3/35-environmental-and-resource-economics/out/scc.csv};
  \legend{1.5~\%, 2.0~\%, 2.5~\%}
\end{axis}
\end{tikzpicture}
\end{omchart}
% figdata: scc
% ledger: b3er:epa

For a tonne emitted in 2020, the estimate is 120 dollars at a near-term rate of 2.5~\%, 190 at 2.0~\% and 340 at 1.5~\%: one point of discounting
nearly triples it. The estimates rise with the year of emission, to 280, 410 and 600 dollars in 2080, because damages
grow with the size of the economy.

\begin{record}[Discount rates in climate assessments]
A widely read review of the economics of climate change concluded that sharp and immediate cuts in emissions
were warranted. Its conclusion depends decisively on a near-zero rate of pure time
preference combined with a particular utility function; with assumptions consistent
with market real interest rates and saving rates it does not hold. % ledger: b3er:nordhaus
Defenders of a low rate reply that the rate of pure time preference is an ethical
choice about the weight of people not yet born, which market rates, set by people
saving for themselves, do not reveal. And another critic of the review's model,
sceptical of its low rate, argued that the case for early action may rest elsewhere:
on insurance against a small chance of ruinous catastrophe, which a calculation built
on expected damages understates. % ledger: b3er:weitzman
The disagreement is not about the climate model but about $\rho$ and $\eta$: the \omterm{def:b3:environmental-and-resource-economics:scc}{social cost of
carbon} is a calculation whose range is mostly a choice of these two numbers, and an
honest estimate states them.
\end{record}

\section{Exhaustible resources}\label{sec:b3:environmental-and-resource-economics:hotelling}

\begin{definition}[Hotelling rule; scarcity rent; backstop technology]\label{def:b3:environmental-and-resource-economics:hotelling}
\index{Hotelling rule}\index{scarcity rent}\index{backstop technology}
For an exhaustible resource extracted at cost $c$ a unit, the \emph{scarcity rent} is the
net price $P_t - c$. The \emph{Hotelling rule} says that in a competitive market with a known
stock it rises at the rate of interest. A \emph{backstop technology} is a substitute
available in unlimited quantity at a constant price, which caps the resource's price.
\end{definition}

\begin{proposition}[The Hotelling rule by arbitrage]\label{prop:b3:environmental-and-resource-economics:hotelling}
If owners can extract a unit now or later and the market is competitive, any path
on which both now and later extraction take place has $P_{t+1} - c = (1 + r)(P_t - c)$. With demand $q(P)$ and a backstop at price $\bar P$, the
initial rent is the one at which cumulative extraction equals the stock exactly when
the price reaches $\bar P$.
\end{proposition}

\begin{proof}
A unit left in the ground is an asset. Extracting it at $t$ yields $P_t - c$, which invested
gives $(1 + r)(P_t - c)$ at $t + 1$; extracting it at $t + 1$ yields $P_{t+1} - c$. If the second exceeded the first,
no owner would extract at $t$; if it fell short, none would wait. Extraction in both
periods requires equality. If the stock ran out before the price reached $\bar P$, the
price would jump at exhaustion, so owners would hold back; if it lasted beyond, the
resource would be unsold at $\bar P$ and the rent would fall. Only the rent that exhausts the
stock at $\bar P$ survives.
\end{proof}

\begin{model}{The Hotelling rule}
\assumes A known stock; competitive owners; extraction cost constant over time.
\predicts Real resource prices, net of extraction cost, rising at the rate of
interest.
\tested Observed prices are often inconsistent with the rule: over long periods
real prices of many minerals have not risen at the rate of interest (Krautkraemer
1998). % ledger: b3er:kra
In 1980 two scholars bet on whether the real price of a bundle of five metals would
rise or fall over the decade; it fell, and the one who bet on falling prices won.
Over the price history of the same bundle from 1900 to 2008, ten-year bets on rising
prices would have won a majority of the time, by a wide margin
(Kiel, Matheson and Golembiewski 2010): the record has no single direction, as the
breaks below suggest. % ledger: b3er:kmg
\breaks Exploration adds to known stocks, ores differ in quality, technical change
lowers extraction costs, and capital in mines makes extraction slow to adjust; all
change the price path. % ledger: b3er:kra
\end{model}

\section{Learning curves in energy}\label{sec:b3:environmental-and-resource-economics:learning}

\begin{definition}[Learning rate]\label{def:b3:environmental-and-resource-economics:learning-rate}
\index{learning rate}
If unit cost falls with cumulative output $Q$ as $C = C_1 Q^{-b}$, the \emph{learning rate} is the
fraction by which cost falls with each doubling of cumulative output, $1 - 2^{-b}$.
\end{definition}

The derivation is one line: $C(2Q)/C(Q) = 2^{-b}$, whatever $Q$. A \omterm{def:b3:environmental-and-resource-economics:learning-rate}{learning rate} of 20~\% means $b = -\log_2 0.8 = 0.32$. Forecasts of
energy costs built on such curves, and tested on more than 50 technologies, imply that most
energy--economy models overestimated the cost of moving away from fossil fuels,
because they underestimated how fast the costs of renewables would fall and how fast
they would be deployed. % ledger: b3er:way
Learning curves are fitted on the past, and nothing guarantees that a rate holds
for another doubling; the forecasts are distributions for that reason.

\section{Exercises}\label{sec:b3:environmental-and-resource-economics:exercises}

\begin{exercise}[$\star$]\label{exo:b3:environmental-and-resource-economics:1}
From the opening chart, give the peak year and level of emissions per person for the
United Kingdom and China, and the incomes at the two peaks.
\end{exercise}

\begin{exercise}[$\star$]\label{exo:b3:environmental-and-resource-economics:2}
Two firms have \omterm{def:b3:environmental-and-resource-economics:abatement}{marginal abatement costs} $2x_1$ and $3x_2$ (dollars a tonne, $x$ in tonnes). Find the cheapest way to cut
100 tonnes, and the tax that achieves it.
\end{exercise}

\begin{exercise}[$\star$]\label{exo:b3:environmental-and-resource-economics:3}
A mineral's net price is 50 dollars a tonne and the interest rate 3~\%. Under the \omterm{def:b3:environmental-and-resource-economics:hotelling}{Hotelling rule}, what is it in ten
years?
\end{exercise}

\begin{exercise}[$\star$]\label{exo:b3:environmental-and-resource-economics:4}
What price data would contradict the \omterm{def:b3:environmental-and-resource-economics:hotelling}{Hotelling rule}, and what would you check
before concluding that they do?
\end{exercise}

\begin{exercise}[$\star\star$]\label{exo:b3:environmental-and-resource-economics:5}
A tonne of CO$_2$ causes 1 dollar of damage a year forever. Find its social cost at constant
discount rates of 1.4, 3 and 5~\%.
\end{exercise}

\begin{exercise}[$\star\star$]\label{exo:b3:environmental-and-resource-economics:6}
Find the Ramsey discount rate for $\rho = 0.1~\%$, $\eta = 1$, $g = 1.3~\%$, and for $\rho = 1.5~\%$, $\eta = 2$, $g = 1.3~\%$.
\end{exercise}

\begin{exercise}[$\star\star$]\label{exo:b3:environmental-and-resource-economics:7}
Find the year, the level and the income per person at the peak of sulphur dioxide
emissions per person in France, Italy and Spain.
\end{exercise}

\begin{exercise}[$\star\star$]\label{exo:b3:environmental-and-resource-economics:8}
Under uncertainty about \omterm{def:b3:environmental-and-resource-economics:abatement}{marginal abatement cost} (a shift of variance $\sigma^2$), a tax beats a cap by
$\sigma^2(s_C - s_D)/(2s_C^2)$, where $s_C$ and $s_D$ are the slopes of \omterm{def:b3:environmental-and-resource-economics:abatement}{marginal abatement cost} and marginal damage. With $\sigma = 10$
and $s_C = 8$, find the advantage when $s_D = 2$ and when $s_D = 20$. Which pollutant is each like?
\end{exercise}

\begin{exercise}[$\star\star$]\label{exo:b3:environmental-and-resource-economics:9}
What data would show that the fall of sulphur dioxide is due to regulation, not to
income itself?
\end{exercise}

\begin{exercise}[$\star\star\star$]\label{exo:b3:environmental-and-resource-economics:10}
Prove the \omterm{def:b3:environmental-and-resource-economics:hotelling}{Hotelling rule} in continuous time: show that an owner indifferent between
extracting at $t$ and at $t + dt$ needs $\dot P/(P - c) = r$, and solve for $P_t$.
\end{exercise}

\begin{exercise}[$\star\star\star$]\label{exo:b3:environmental-and-resource-economics:11}
A technology's unit cost fell from 4 to 0.4 dollars while cumulative output grew a thousandfold. Find $b$ and the
\omterm{def:b3:environmental-and-resource-economics:learning-rate}{learning rate}.
\end{exercise}

\begin{exercise}[$\star\star\star$]\label{exo:b3:environmental-and-resource-economics:12}
A tonne emitted today causes 1.5 dollars of damage in a year's time, growing 2~\% a year, for 300 years. Find its
present value at constant rates of 1.4, 3 and 5~\%, and with a rate that falls from 3 to 1.5~\% over the first
century.
\end{exercise}

\section{Problem: Brenland's Coal and Its Smoke}\label{sec:b3:environmental-and-resource-economics:problem}

\begin{problem}[{Weekend problem --- Brenland's coal and its smoke}]\label{pb:b3:environmental-and-resource-economics:coal}
An invented Brenland has 600 million tonnes of coal in 2025, extracted at 20 dalens a tonne. Demand is
$q = 40 - 0.25P$ million tonnes a year; a backstop supplies the same energy at a coal-equivalent 100 dalens; the
interest rate is 4~\%. Each tonne of coal burnt emits 0.02 tonnes of sulphur dioxide; each kilotonne of it a
year causes 0.4 deaths; the \omterm{def:b3:cost-benefit-analysis:vsl}{value of a statistical life} is 3.0 million dalens. Two power firms abate with marginal
costs $10x_1$ and $30x_2$ dalens a tonne ($x$ in kilotonnes). Brenland emits 40 million tonnes of CO$_2$ a year. % figdata: coal_stats

\medskip\noindent\textbf{Part I --- the Hotelling path.}
\begin{enumerate}
\item Write the price path and the condition that fixes the initial rent.
\item Find the rent and the price in 2025.
\item Find extraction in 2025.
\item Find the rent, the price and extraction in 2035.
\item When does the coal run out?
\item What would a discovery of more coal do to the rent and the date?
\item What would a cheaper backstop do?
\end{enumerate}

\medskip\noindent\textbf{Part II --- the smoke.}
\begin{enumerate}[resume]
\item Find sulphur dioxide emissions in 2025.
\item Find the deaths they cause each year.
\item Find the marginal damage of a tonne.
\item Find the total yearly damage.
\item Which number in the chain is the most uncertain?
\item Why should the Hotelling path change once the smoke is priced?
\end{enumerate}

\medskip\noindent\textbf{Part III --- tax, permits or standards.}
\begin{enumerate}[resume]
\item Find the efficient tax on sulphur dioxide.
\item Find each firm's abatement under it.
\item Find the total abatement cost.
\item Find the cost of the same total cut split equally, and the saving from the tax.
\item What cap and permit price give the same result?
\item Find the deaths avoided and the money value of the abatement's benefit.
\item Why might a regulator still choose uniform standards?
\end{enumerate}

\medskip\noindent\textbf{Part IV --- carbon.} Value Brenland's emissions with the EPA's 2023 figures for 2020
emissions, read as dalens per tonne.
\begin{enumerate}[resume]
\item Find the yearly social cost at a near-term rate of 2~\%.
\item Find it at 2.5 and 1.5~\%.
\item Express the 2~\% figure as a share of Brenland's 2025 GDP (\code{load("accounts")}).
\item With $g = 1.3~\%$, give two pairs $(\rho, \eta)$ that give about 2~\%.
\item The \omterm{def:b3:environmental-and-resource-economics:scc}{social cost of carbon} counts damages everywhere. Why does that matter for
the tax a single country would choose for itself?
\item State the result: the efficient sulphur dioxide tax and the year Brenland's coal is
exhausted.
\end{enumerate}
\end{problem}
